\begin{abstract}
随着云计算、移动互联网、人工智能和物联网技术的快速发展，端到端网络传输正在从单一固定终端和局域高质量链路，扩展到跨地域远距离传输、蜂窝/无线/卫星等多接入网络，以及手机、笔记本、台式机和边缘设备等多模态终端并存的复杂场景。传统基于丢包的拥塞控制算法（如TCP NewReno、CUBIC）存在拥塞信号滞后、高带宽延迟积（Bandwidth-Delay Product，BDP）网络收敛缓慢、无法区分拥塞程度等固有缺陷；而基于延迟的算法（如TCP Vegas、BBR）则面临往返时延（Round-Trip Time，RTT）测量不稳定、与丢包算法共存时公平性差等问题。本文提出Swift算法——一种启发式多信号融合自适应网络拥塞控制方法。该算法名称寓意"雨燕"，象征着高吞吐量（High Throughput）与低延迟（Low Latency）的设计目标，如同雨燕般轻盈敏捷、反应迅速。Swift的决策以确定性规则在每个确认事件上于协议栈之外求解，不依赖训练样本与神经网络推理，行为可逐分支解释、结果可完全复现。

Swift算法的核心创新收敛为三点：（1）面向远距离与异构终端传输的11维有效观测子空间设计，该子空间嵌入于15元素状态传输容器，另含4项跨进程路由用的元数据通道，将显式拥塞通知（Explicit Congestion Notification，ECN）状态、拥塞控制状态机（Congestion Algorithm State）状态、拥塞事件类型、RTT和在途字节等协议栈信息纳入决策状态表示；（2）多信号融合拥塞判定与差异化响应机制，在协议栈的显式降窗回调内部按拥塞状态机状态与ECN状态完成"超时/ECN驱动的窗口缩减/普通丢包"三路语义判别，并针对三类信号分别采用0.50、0.75、0.70的窗口保留因子（Window Retention Factor），同时引入最多连续3次缩减后保持窗口的连续缩减保护与降窗后4个ACK事件冻结机制，以区分拥塞信号的性质和严重程度并抑制窗口快速反弹；（3）启发式反馈自适应的融合控制策略，根据最小RTT感知的RTT膨胀阈值、性能反馈值快慢双指数移动平均（Exponential Moving Average，EMA）的基线相对偏离以及连续增长趋势动态调整乘性增加因子$\alpha$（取值范围$[0.85, 1.30]$，基准值1.10），并以"时间窗口投递速率采样加窗口化最大值过滤（Windowed Max-filter）"两级结构估计带宽延迟积（Bandwidth-Delay Product，BDP），在拥塞避免阶段围绕$\alpha \times BDP$目标窗口进行有界步长的渐进拉升或回落。

本文基于ns-3.40网络仿真平台实现了Swift算法原型，仿真进程与决策进程经ZeroMQ同步请求—应答信道交互（由ns3-gym/OpenGym框架承载15元素状态容器的跨进程传输，仅作为消息通道使用，不含学习逻辑）。实验矩阵覆盖近端高速、共享收敛链路、无线局域网接入、蜂窝移动接入、城域与长途/跨域广域网及LEO/GEO卫星通信共36个场景，并通过\texttt{enable\_udp\_burst}开关构建"纯TCP"与"TCP $+$ 平均负载为瓶颈速率32\%（峰值64\%，50\%占空比）的UDP Burst"两类A/B严格对照，共288次仿真。全部指标由ns-3 FlowMonitor产物重新导出并\emph{仅在前向数据流上定义}；瓶颈队列统一采用启用ECN标记的RED（$MinTh=0.3Q$、$MaxTh=0.9Q$），四种协议对称开启ECN协商。288条记录全部通过物理相容性与取值域审计，无需排除数据点。\textbf{稳态实验结果}表明：纯TCP设置下Swift在\textbf{全部36个场景}取得最高聚合吞吐量，相对最强基线$+2.1\%\sim+5.8\%$（均值$+4.1\%$），且增益在近端高速、无线/蜂窝接入与远距离链路三类路径上均匀分布；平均单向时延相对NewReno/CUBIC均值平均低约2.3\%（36个场景中24个更低），与BBR基本相当；四种协议全部场景丢包率不超过0.026\%，Swift的Jain公平性指数均值0.9989为四者最高。\textbf{跨流量鲁棒性实验}表明：突发下Swift的绝对吞吐量仍在36/36个场景领先（相对最强基线$+2.2\%\sim+5.9\%$），吞吐量保持率均值68.6\%与基线持平、无任何协议发生塌缩；GEO卫星场景纯TCP吞吐量8.78\,Mbps（相对NewReno $+7.5\%$）、突发下5.65\,Mbps（相对NewReno $+10.1\%$），LEO卫星与长途WAN场景同样为四协议最高。本文同时如实披露混合结果与边界：Swift的丢包率与NewReno相当、略低于CUBIC/BBR（差值不超过0.003个百分点），本文不主张其"更不易丢包"；个别深缓冲场景中Swift相对延迟最优基线的时延抬升最高约10\%；每(场景,协议,设置)为单一随机种子的确定性运行，结论以36/36的正增益方向一致性为主要支撑。相应的改进方向与证据范围限制（单一随机种子、RED/ECN配置适用范围）已在第5章给出。

\keywords{拥塞控制，启发式自适应，多信号融合，显式拥塞通知，自适应参数调优，远距离传输，异构终端}
\end{abstract}

\begin{englishabstract}

With the rapid development of cloud computing, mobile Internet, artificial intelligence, and Internet-of-Things technologies, end-to-end network transmission is expanding from fixed terminals and local high-quality links to long-distance cross-region paths, cellular/wireless/satellite access networks, and heterogeneous endpoint environments involving phones, laptops, desktops, edge devices, and other terminal classes. Traditional loss-based congestion control algorithms (such as TCP NewReno and CUBIC) suffer from delayed congestion signals, slow convergence in high Bandwidth-Delay Product (BDP) networks, and inability to distinguish congestion severity. Meanwhile, delay-based algorithms (such as TCP Vegas and BBR) face unstable Round-Trip Time (RTT) measurements and poor fairness when coexisting with loss-based algorithms. This paper proposes Swift, a \emph{heuristic} multi-signal fusion adaptive network congestion control algorithm. The name ``Swift'' symbolizes the design goals of High Throughput and Low Latency, like a swift that is light, agile, and responsive. All of Swift's decisions are computed by deterministic rules outside the protocol stack on every acknowledgement event, without any training samples or neural-network inference, so its behaviour is branch-wise explainable and fully reproducible.

The core innovations of Swift are consolidated into three aspects: (1) an 11-dimensional effective observation subspace for long-distance and heterogeneous endpoint transmission, embedded in a 15-element state transport container with four additional metadata channels for inter-process routing, which incorporates Explicit Congestion Notification (ECN) status, Congestion Algorithm (CA) state, congestion event types, RTT, and bytes in flight into the decision state representation; (2) a multi-signal congestion determination and differentiated response mechanism that, inside the protocol stack's explicit window-reduction callback, performs a three-way semantic classification into timeout, ECN-driven window reduction, and ordinary loss according to the CA state and the ECN state, applying window retention factors of 0.50, 0.75, and 0.70 respectively, together with a consecutive-decrease guard that holds the window after at most three successive reductions and a four-ACK post-decrease freeze; and (3) a heuristic feedback-adaptive fusion control strategy that adjusts the multiplicative increase factor $\alpha \in [0.85, 1.30]$ (baseline 1.10) from min-RTT-aware RTT-inflation thresholds, the \emph{baseline-relative} deviation between a fast and a slow performance-feedback Exponential Moving Average (EMA), and consecutive growth trends, and that estimates the Bandwidth-Delay Product (BDP) through a two-stage structure combining time-window delivery-rate sampling with a windowed max-filter, moving the congestion window toward the $\alpha \times BDP$ target with bounded pull-up or drain steps during congestion avoidance.

We implemented the Swift prototype on the ns-3.40 simulation platform; the simulator and decision processes interact over a synchronous ZeroMQ request--reply channel (the 15-element state container is carried by the ns3-gym/OpenGym framework, used purely as a message channel and containing no learning logic). The experiment matrix spans 36 scenarios covering near-end high-speed links, shared oversubscribed links, wireless LAN access, cellular mobile access, metropolitan/long-haul/inter-region wide area networks, and LEO/GEO satellite links, and an \texttt{enable\_udp\_burst} switch constructs a strict A/B comparison between ``pure TCP'' and ``TCP $+$ an On-Off UDP burst whose average load equals 32\% of the bottleneck rate (64\% peak, 50\% duty cycle)'', for 288 simulation runs in total. All metrics were re-derived from the ns-3 FlowMonitor artifacts and are defined \emph{on the forward data flows only}; the bottleneck queue is a RED queue with ECN marking ($MinTh=0.3Q$, $MaxTh=0.9Q$) and ECN negotiation is enabled symmetrically for all four protocols. All 288 records passed integrity, physical-consistency, and value-domain audits, and no data point needed to be excluded. \textbf{Steady-state results} show that, under the pure-TCP setting, Swift attains the highest aggregate throughput in \emph{all 36 scenarios}, with a gain of $+2.1\%$--$+5.8\%$ (mean $+4.1\%$) over the strongest baseline; the gain is distributed uniformly across near-end high-speed links, wireless/cellular access, and long-distance paths. Its mean one-way delay is on average about 2.3\% lower than the NewReno/CUBIC mean (lower in 24 of 36 scenarios) and comparable to BBR; the loss rates of all four protocols stay below 0.026\% in every scenario, and Swift's mean Jain fairness index of 0.9989 is the highest. \textbf{Cross-traffic robustness experiments} further show that Swift keeps the highest absolute throughput in 36/36 paired scenarios under the burst ($+2.2\%$--$+5.9\%$ versus the strongest baseline), with a mean throughput-retention ratio of 68.6\% on a par with the baselines and no collapse for any protocol. In the GEO satellite scenario Swift reaches 8.78\,Mbps under pure TCP ($+7.5\%$ versus NewReno) and 5.65\,Mbps under the burst ($+10.1\%$ versus NewReno); it also leads in the LEO satellite and long-haul WAN scenarios. Mixed results and boundaries are reported faithfully: Swift's loss rate is comparable to NewReno and slightly lower than CUBIC/BBR (by at most 0.003 percentage points), and this paper does not claim that Swift is ``less loss-prone''; in a few deep-buffer scenarios Swift's delay can be up to about 10\% above the lowest-delay baseline; each (scenario, protocol, setting) triple corresponds to a single deterministic run, and the conclusions are primarily supported by the consistency of the 36/36 positive-gain direction. Refinement directions and evidence-scope limits (single random seed, RED/ECN configuration scope) are discussed in Chapter 5.

\englishkeywords{Congestion Control, Heuristic Adaptation, Multi-signal Fusion, Explicit Congestion Notification, Adaptive Parameter Tuning, Long-distance Transmission, Heterogeneous Endpoints}

\end{englishabstract}
